% Sections 5.1-5.5, from lectures/L05/appendix/a_variables_and_hardware.md (A.1-A.4).

\section{What is left}
\label{c5:sec:intro}

Every construct you need in order to describe hardware you have already used: \code{entity} and
\code{architecture} in \chapref{c1:ch}, \code{std_logic_vector}, \code{process} and \code{case} in
\chapref{c2:ch}, the clocked process and the scheduling rule in \chapref{c3:ch}, subcomponents and
generics in \chapref{c2:ch} and \ref{c4:ch}.

Two things are left, and they are the two hardest to work out on your own:
\begin{itemize}
  \item The first is \code{variable}, the one construct the book has not needed until now. It looks
    like a signal and behaves nothing like one, and the difference is exactly the kind that gives
    you a design which simulates plausibly and is silently wrong. \Secref{c5:sec:variable} works
    through the case where it bites.
  \item The second is what any of this actually becomes. Since \chapref{c1:ch} you have trusted the
    synthesis tool to build the circuit you described. It does, but not out of anything you would
    recognise. \Secref{c5:sec:fabric} opens the box. It turns out an FPGA contains no gates at all,
    which is why \chapref{c2:ch}'s Karnaugh maps bought you less than they appeared to. Your clock
    has a speed limit (\secref{c5:sec:fmax}), and something specific sets it. And a \code{case} with
    a missing branch silently gives you a memory element nobody asked for
    (\secref{c5:sec:latch}).
\end{itemize}

Neither section answers ``is my logic right?'' You have been answering that question since
\chapref{c1:ch}. They answer the two questions after it: is this what I meant, and will it fit and
run.

% ------------------------------------------------------------------------------------------
\section{\code{variable}: process-local storage that updates immediately}
\label{c5:sec:variable}

A \code{variable} is declared inside a \code{process}, between \code{process(...)} and
\code{begin}. The decisive difference from a \code{signal}: variable assignment uses \code{:=} and
takes effect \textbf{immediately}, so the very next statement that reads it sees the updated value,
exactly like an assignment in C. There is no scheduling, none of the
wait-until-the-process-suspends delay \secref{c3:sec:template} spelled out for signals, and no
``last assignment wins''.

\begin{booktable}{@{}p{9em}p{9.6em}L@{}}
   & \thead{\code{signal} (\code{<=})} & \thead{\code{variable} (\code{:=})} \\
  \midrule
  Declared & In the architecture's declarative part & In the process's declarative part \\
  Visible from & Anywhere in the architecture & Only inside its own process \\
  Update time & Scheduled: applied when the process next suspends & Immediate: applied before the
    next statement \\
  Repeated assignment in one round & Only the last survives; earlier ones are discarded & Each one
    takes effect, in turn \\
\end{booktable}

\subsection*{Why this actually matters: an XOR parity generator}
An entity computing the XOR parity of an 8-bit input, \code{'1'} if \code{bits} has an odd number
of bits set:

\begin{vhdlcode}
entity parity_gen is
    port(bits  : in  std_logic_vector(7 downto 0);
         parity: out std_logic);
end entity;
\end{vhdlcode}

Coming from software, the first instinct is to accumulate the running XOR in a loop, with a signal
as the accumulator:

\begin{vhdlcode}
-- Don't do this: shown to demonstrate why it doesn't work.
architecture broken of parity_gen is
signal parity_s: std_logic;
begin
    parity <= parity_s;

    process(bits) is
    begin
        parity_s <= '0';
        for i in bits'range loop
            parity_s <= parity_s xor bits(i);
        end loop;
    end process;
end architecture;
\end{vhdlcode}

Trace one round, with the scheduling rule from \secref{c3:sec:template}: every read of
\code{parity_s} returns what it held \emph{before this round began} (call it \code{old_parity}),
and only the last assignment survives.
\begin{itemize}
  \item \code{parity_s <= '0';} schedules \code{'0'}. Not applied yet.
  \item \code{i = 7} (\code{bits'range} on \code{(7 downto 0)} runs from high to low): reads
    \code{parity_s}, gets \code{old_parity}, since the \code{'0'} has not been applied. Schedules
    \code{old_parity xor bits(7)}, discarding the \code{'0'}.
  \item \code{i = 6}: reads \code{parity_s}, again \code{old_parity}. Schedules
    \code{old_parity xor bits(6)}, discarding the previous one.
  \item \dots and so on for every iteration \dots
  \item \code{i = 0}, the last statement of the round: still reads \code{old_parity}, schedules
    \code{old_parity xor bits(0)}, discarding everything before it.
  \item The process suspends and \code{parity_s} becomes \code{old_parity xor bits(0)}, the only one
    of nine scheduled assignments that survived.
\end{itemize}

The result depends only on \code{bits(0)} and on whatever \code{parity_s} happened to hold already.
Seven iterations and the initial \code{'0'} were computed and thrown away. This is no edge case: it
is what a signal-as-accumulator \emph{always} does, because every iteration reads the same frozen
value and only the textually last assignment is kept.

The same computation with a \code{variable}, and this is the worked example on disk:

\vhdlfile{lectures/L05/parity_gen/parity_gen.vhd}

Now \code{acc := '0'} is applied immediately; \code{i = 7} reads \code{'0'}, XORs in \code{bits(7)}
and updates \code{acc} at once; \code{i = 6} reads the just-updated value, and so on through every
iteration. The loop ends with \code{acc} holding the true XOR of all eight bits, and
\code{parity_s <= acc;} schedules the single correct value, the only assignment to \code{parity_s}
in the round.

\subsection*{What that loop is not}
It is not eight steps the hardware performs one after another. This is where a background in C
misleads most, so it is worth saying outright: a \code{for} loop whose bounds are fixed at
elaboration is \textbf{unrolled} by the synthesis tool. The eight iterations become eight XOR gates
in a chain, all settling within one propagation delay, in the same clock cycle. There is no counter
for \code{i}, no branching, no iteration in time. The loop writes eight gates in three lines; it
does not do eight things in sequence.

That is also why \code{acc}, which reads like a running total, costs no storage at all. It never
has to survive a clock edge: each of its values is just the wire between two of those XOR gates. A
\code{variable} is the right tool precisely because it never survives a round.

The same goes for \code{bits'range}. The \textbf{attribute} \code{'range} gives the index range of
\code{bits}, here \code{7 downto 0}, so that the loop covers every bit without the bounds being
written twice. It is resolved at elaboration, like everything else in the loop's header, which is
exactly why the tool can unroll it. \code{for i in 0 to 1 loop} in \chapref{c3:ch}'s
\file{led_toggle.vhd} means the same thing: both bits get gates of their own and both update in the
same cycle.

\subsection*{The rule of thumb}
\begin{itemize}
  \item Use a \code{variable} for a temporary computation consumed entirely within one round: loop
    accumulators, temporary swaps, intermediate values handed over to a signal at the end, exactly
    like \code{acc}.
  \item Use a \code{signal} for anything another process or the outside world reads, and anything
    that has to hold its value across separate rounds (a counter, or a shift register's stored
    bits, both in \chapref{c6:ch}).
\end{itemize}

A variable \emph{can} keep its value between rounds through the same process, behaving as static
local storage rather than a fresh local variable each time, so it is possible to build state with
one. But since its updates are immediate rather than scheduled, it behaves differently from a signal
used the same way, particularly for anything modelling a chain of registers stepping one place per
clock edge. If state genuinely has to persist and be read by other logic, reach for a
\code{signal}: the scheduled-update semantics are what make \chapref{c6:ch}'s counters and shift
registers work.

% ------------------------------------------------------------------------------------------
\section{What the FPGA actually consists of}
\label{c5:sec:fabric}

An FPGA is not a blank slate that gates are etched into. It is a fixed grid of identical small
blocks, manufactured before anyone wrote your VHDL, and ``synthesis'' is working out how to
configure them so that they behave like the circuit you described. Two of them matter here:
\begin{itemize}
  \item A \textbf{lookup table}, or \textbf{LUT}: a small memory, typically four to six inputs and
    one output. Load it with the truth table of any Boolean function of those inputs and it
    computes that function. That is the honest answer to ``what does a gate become on an FPGA'':
    nothing becomes a gate. A four-input LUT computes \code{a and b},
    \code{a xor b or (c and not d)}, or any other function of up to four inputs at exactly the same
    cost, since in every case it is the same block holding a different table.
  \item A \textbf{flip-flop}: one bit of clocked storage, sitting beside every LUT.
\end{itemize}

On the DE0-CV's Cyclone~V they are packaged together as an \textbf{ALM}, which holds a splittable
six-input LUT, usable as one function of six variables or two smaller ones, plus four flip-flops.
That is the unit Quartus's fitter report counts, which is why it talks about ALMs rather than
gates.

That is almost the whole story. \code{register4} becomes four flip-flops, \code{mux_8to1} a handful
of LUTs, and \chapref{c8:ch}'s state machine LUTs feeding the flip-flops holding the state.

It also explains something from \chapref{c2:ch} that would otherwise look like a broken promise:
the book had you minimise expressions with Karnaugh maps to save gates, and then never mentioned the
saving again. On an FPGA a function of four variables costs one LUT whether you minimised it or not.
Minimisation still matters: it is how you understand a circuit, it is exactly right for the discrete
logic and ASIC flows Karnaugh maps were invented for, and reducing a function below the LUT's input
width is what keeps it from needing a second level of LUTs. But on this board it usually buys
nothing, and it is better to know that than to wonder.

% ------------------------------------------------------------------------------------------
\section{Why a clock has a speed limit}
\label{c5:sec:fmax}

Signals take time through a LUT and along the wires between blocks. That \textbf{propagation delay}
is why a clock cannot run arbitrarily fast.

Picture the shape every design in this book has, a flip-flop, some combinational logic, another
flip-flop:
\begin{itemize}
  \item At a rising edge the first presents a new value, which has to travel through every LUT and
    wire between the two and \emph{settle} before the next edge arrives.
  \item If it has not, the second flip-flop captures a value that is still changing, and the design
    does the wrong thing reliably and mysteriously.
\end{itemize}

So the clock period has to cover three things:
\begin{itemize}
  \item The first flip-flop's \textbf{clock-to-Q} delay: the time between the edge and its output
    actually changing.
  \item The \textbf{combinational delay} through the LUTs and wires between the two.
  \item The second flip-flop's \textbf{setup time}.
\end{itemize}

The slowest such path is the \textbf{critical path}, and one over that sum is \textbf{Fmax}.

Quartus computes Fmax for you, in the timing analysis step (appendix~\ref{app:quartus}), but only
if you tell it what the clock is. That means an \file{.sdc} constraints file with a
\code{create_clock} on the \code{50 MHz} input. Without one the analyser has no period to compare
anything against, so it reports every path as unconstrained and gives you no Fmax at all. Distinct
from that, a design that ``misses timing'' is one where some path turned out slower than the period
you actually declared.

\textbf{None of this is part of this book.} You will not write an \file{.sdc}, read a timing report,
or meet a design that misses timing, since every design here fits inside 20~ns with room to spare.
It is worth knowing that the step exists and what it requires, so that you recognise it the first
time a design of yours approaches the limit.

Two consequences worth carrying:
\begin{itemize}
  \item \textbf{Deep combinational logic costs you speed, not the number of flip-flops.} Chaining a
    long expression between two registers lengthens the critical path; splitting it so that part of
    the work happens on one cycle and the rest on the next shortens it. That trade-off, more
    registers against a faster clock, is the single most common move in fast digital design.
  \item \textbf{\code{50 MHz} is a budget of 20~ns.} That is the number every path in your design
    competes against, and it is plenty of time at these speeds.
\end{itemize}

% ------------------------------------------------------------------------------------------
\section{The latch you get by accident}
\label{c5:sec:latch}

\Chapref{c3:ch} built a D latch, showed that it was transparent while enabled, and then set it
aside in favour of the flip-flop. Here it comes back uninvited.

\Secref{c2:sec:process} required a \code{when others} on every \code{case}, and gave the language's
reason: a \code{case} must cover every value of its selector's type. The general rule behind it is
that a combinational \code{process} must assign its output on \emph{every} path through it, and
that reason can now be stated:
\begin{itemize}
  \item VHDL requires a signal to keep its old value if nothing assigns a new one.
  \item If some path through your process leaves \code{x} unassigned, you have said that \code{x}
    has to \emph{remember} its previous value on that path, and remembering is not something a
    combinational network can do.
  \item So the tool builds the only thing that can: a latch, exactly the one from
    \secref{c3:sec:latch}.
\end{itemize}

That is an \textbf{unintended latch}, and it is almost never what anyone meant. It is a warning
rather than an error, the design will often appear to work, and the message is worth recognising:

\begin{console}
Warning: Inferred latch(es) for signal "x"
\end{console}

The fix is never to add a latch deliberately. It is to assign the signal on every path, which is the
rule you were already following.
